% !TeX root = ../../main.tex
\subsection{Exact low-genus normalization checks}
\label{subsec:c:01}

We use the exact vertices \eqref{eq:s06:exact-vertex}, derived from
Faber--Pandharipande's formula \cite[Theorem 1]{FP03}, and the
complete tree identity \eqref{eq:s06:tree-expansion}. We first
classify the small trees, then integrate their vertex polynomials
with their automorphism weights. Mumford's relation
\cite[\S 5]{Mum83} supplies the remaining normalizations.

\begin{prp}[Low-genus values]\label{prp:c:low-genus}
The following values are exact:
\begin{center}
\begin{tabular}{c|c|c|c}
$g$ & $a_g$ & $b_g$ & $(b_g(h)/b_g)_{1\leq h\leq g}$\\
\hline
$2$ & $1/2880$ & $1/5760$ & $(2,2)$\\
$3$ & $1/90720$ & $1/725760$ & $(4,8,8)$
\end{tabular}
\end{center}
The genus-three fixed-insertion normalization is
\begin{equation}\label{eq:c:lambda-three}
 \langle\lambda_3\rangle_3=\frac18.
\end{equation}
\end{prp}
\begin{prf}
\emph{The complete tree lists.}
For a tree with $E$ edges, the vertex orders
$m_v=2g_v-3+n_v$ are nonnegative and sum to $g_2-E$,
by \eqref{eq:s06:tree-degree}. Thus $E\leq1$ in genus two
and $E\leq3$ in genus three. A leaf must have positive genus.
With one edge this gives two elliptic vertices in genus two,
and genera one and two in genus three. With two edges the tree
is a chain; its bivalent middle vertex also has positive genus,
so genus three forces all three genera to be one. With three
edges in genus three, every $m_v$ is zero. There is no bivalent
vertex of order zero, so the tree has a trivalent rational center
and three elliptic leaves. Including the one-vertex trees,
these lists are complete.

\smallskip\noindent
\emph{Exact vertices.}
For edge parameters $s,t,u\in[0,1]$, finite coefficient extraction
in \eqref{eq:s06:exact-vertex} gives
\begin{equation}\label{eq:c:vertices}
 \begin{aligned}
 W^{\rm top}_{2,0}&=\beta_2,&
 W^{\rm top}_{3,0}&=\frac{61}{6}\beta_3,\\
 W^{\rm top}_{1,1}(s)&=\beta_1,&
 W^{\rm top}_{0,3}(s,t,u)&=1,\\
 W^{\rm top}_{2,1}(s)&=\beta_2(5/2-3s+s^2/2),\\
 W^{\rm top}_{1,2}(s,t)&=\beta_1(1-s-t).
 \end{aligned}
\end{equation}
Indeed, it suffices to use
$\Phi(x)=1-x/2+x^2/12-x^3/144+O(x^4)$ and
$J_0(2\sqrt{sx})=1-sx+s^2x^2/4-s^3x^3/36+O(x^4)$,
where $x$ is the coefficient variable. The reciprocal sine
series \eqref{eq:s06:beta-series} similarly gives
\[
 \beta_1=\frac1{24},\qquad
 \beta_2=\frac7{5760},\qquad
 \beta_3=\frac{31}{967680}.
\]

\smallskip\noindent
\emph{Edge integrals and weights.}
The elliptic pair has automorphism order two and contributes
$\frac12\int_0^1\beta_1^2\,\dd s=\beta_1^2/2$. The genus-three
one-vertex, mixed-genus edge, elliptic chain and rational star
have orders $1,1,2,6$, respectively. The nonconstant edge
integrals are
\[
 \int_0^1(5/2-3s+s^2/2)\,\dd s=\frac76,\qquad
 \int_0^1\int_0^1(1-s-t)\,\dd s\,\dd t=0.
\]
Thus the elliptic chain cancels after integration; it is not
discarded from the tree list. The star has constant integrand
$\beta_1^3$, and $\int_{[0,1]^3}1\,\dd s\,\dd t\,\dd u=1$.
The genus-three weighted sum is therefore
$61\beta_3/6+7\beta_1\beta_2/6+\beta_1^3/6$.
Restoring the factorials in \eqref{eq:s06:tree-expansion}
multiplies the genus-two sum by $1!$ and the genus-three sum by
$3!$, giving
\begin{equation}\label{eq:c:tree-values}
 \begin{aligned}
 12b_2&=\beta_2+\frac{\beta_1^2}{2},\\
 12^3b_3&=61\beta_3+7\beta_1\beta_2+\beta_1^3.
 \end{aligned}
\end{equation}
Substituting the exact beta values proves the displayed $b_g$ values.

\smallskip\noindent
\emph{Chern-class normalizations.}
Equation~\eqref{eq:s02:chern-exponential} implies
$c_t(\EE)c_{-t}(\EE)=1$. Its degree-two coefficient gives
$\lambda_1^2=2\lambda_2$; in rank three its degree-four coefficient
gives $\lambda_2^2=2\lambda_1\lambda_3$. Consequently
$a_2=2b_2$ and $a_3=8b_3$.
The defect ratios follow from these relations and the definitions:
in particular $b_3(1)=a_3/2$, while the remaining numerators
$b_g(h)$ in the table equal $a_g$. Finally
$\langle\lambda_3\rangle_3=b_3/a_3=1/8$, as asserted.
\end{prf}
